\section{Márgenes, soporte y radio complejo}
Identificamos $z=jD$ con un núcleo entre dos espacios de $j$ puntos de masa $1/j$.
Al anular los $v$ sitios del exterior se pierde $O(M_0v/j)$ como máximo en un margen.
Al suprimir $h$ índices por lado, la medida uniforme en los supervivientes, extendida
por cero al conjunto original, cambia en norma de variación total a lo sumo $2h/j$.
\paragraph{Balanceo.}
Sean $\mu,\nu$ medidas complejas de masa uno cercanas a las originales, y $z'$ el
núcleo con los ceros adicionales. Buscamos $z'e^{s_r+t_c}$ con márgenes uno.
Fijemos la media original de $t$ en cero. Imponemos todas las ecuaciones de fila y
proyectamos las de columna sobre el complemento de las constantes por la media original.
Si los residuos de columna son constantes y todos los de fila se anulan, la constante
vale cero integrando los dos sistemas con $\nu,\mu$, respectivamente.
En los datos iniciales el jacobiano aplicado a $(s,t)$ es $(s+Dt,(I-\Pi)(t+D^ts))$,
con los productos escritos en notación matricial ordinaria y $\Pi=J_j$.
El complemento de Schur en columnas centradas es $I-D^tD$. Su inverso es uniforme en
norma dos normalizada, por $r<1$. Las ecuaciones de cada coordenada dan a continuación
una cota en supremo por el segundo miembro en supremo, por Cauchy--Schwarz y $|z|\le M_0$.
El residuo y la variación del jacobiano al origen cuestan
$O(v/j+\|\mu-\mu_0\|_{\rm TV}+\|\nu-\nu_0\|_{\rm TV})$ en supremo.
En efecto, el núcleo quitado tiene masa pequeña en cada fila y cada columna originales;
contra el cambio de medida se usa su supremo acotado. Las segundas derivadas están
uniformemente acotadas en una bola de potenciales de radio fijo. Aplicar el inverso
inicial al residuo define una iteración contractiva en una bola de radio proporcional
a la perturbación, para ésta pequeña. La solución es holomorfa por convergencia
localmente uniforme. Conservamos las ecuaciones auxiliares de los índices de masa cero.
El procedimiento mantiene los ceros estructurales; no divide por ninguna entrada positiva.
\paragraph{Difusión y brecha físicas de los menores.}
Normalizando con las nuevas medidas obtenemos $mD'\le C$. Para la brecha, en la
representación euclídea incrustamos el núcleo como
$\operatorname{diag}(\sqrt\mu)\,z'e^{s+t}\operatorname{diag}(\sqrt\nu)$,
con producto por entradas dentro del núcleo. El cambio por pesos de medida es
$O(\sqrt{h/j})$ en Hilbert--Schmidt, pues
$\|\sqrt\mu-\sqrt{\mu_0}\|_2^2\le\|\mu-\mu_0\|_1$ y el núcleo es acotado.
Los nuevos ceros, evaluados primero en las medidas iniciales, cuestan
$\le M_0\sqrt v/j$, y el reequilibrado cuesta $O((v+h)/j)$.
La parte constante $\sqrt\mu\,\sqrt\nu^t$ cambia también $O(\sqrt{h/j})$.
Tomando primero $a$ pequeño respecto a la brecha y luego $v/j$ pequeño se obtiene
un $r'<1$ común. Esto cubre cada menor y, en particular, todos los subcasos de una cara.
\paragraph{Núcleos complejos.}
Recentrar alrededor del exterior real ya balanceado permite dejar fijo su soporte.
Sea $O$ el operador entre ambos lados con los dos bloques cruzados $z-1,(z-1)^t$.
En una región $|y w|\le r_1(1+d_0)$ con $r_1>1$ y $r_1(1+d_0)r'<1$,
$(I+yw O)^{-1}$ es acotado también en supremo: en el centro real se usa
$I-ywO+(yw)^2O(I+ywO)^{-1}O$, que regulariza a ambos lados del inverso en norma dos.
Una perturbación pequeña en TV y supremo del núcleo da una perturbación pequeña en
norma de operador sobre el supremo, en la misma representación por funciones.
La serie de Neumann conserva el inverso. El propagador $ywO(I+ywO)^{-1}$
tiene núcleo puntualmente acotado, aplicando la inversa a cada columna acotada.
En complejo las transpuestas se interpretan bilinealmente. Para el determinante de
Fredholm, el término de grado $l$ está mayorado por $C^l l^{l/2}/l!$ por Hadamard.
Telescopar factores de medida y filas del determinante da continuidad uniforme
al variar la medida en TV y el núcleo en supremo. En el centro real se tiene separación
de cero uniforme por la brecha y la suma de cuadrados de los valores singulares.
Con radio de datos menor se mantienen esa separación y una rama acotada del logaritmo
por continuación desde el centro. Esta extensión sólo se usa en los aproximantes.
